\section{Hinweise zu Rechtschreibung und Zeichensetzung}

In der Vergangenheit zeigten sich bei den eingereichten Arbeiten häufig Fehler in einzelnen Bereichen, auf die hier kurz eingegangen werden soll, um diese Fehler möglichst im Voraus zu verhindern.

\subsection{Zusammen- und Getrenntschreibung}

Ein sehr häufiger Fehler ist, dass zusammengesetzte Wörter im Deutschen, sog. „Komposita“, fälschlicherweise „auseinander“, also in einzelnen Wörtern geschrieben werden. Die Grundregel ist jedoch, ein zusammengesetztes Hauptwort, welches aus einzelnen Wörtern besteht, zusammen zu schreiben, also z. B. „Spülmittelkonzentratzusammensetzungsformel“, auch wenn dieses Bandwurmwort sicherlich ein Extremfall ist (und fälschlicherweise(!) von der Rechtschreibprüfung von Word als falsch markiert wird).

Diese Thematik zeigt sich aufgrund verschiedener Einflüsse als besonders problematisch, da jeder im Umfeld häufig mit falschem Schreiben konfrontiert wird, insbesondere:

\begin{itemize}
    \item durch den Einfluss der englischen Sprache, wo häufig Dinge „auseinander geschrieben“ werden, wie z.\,B. „detergent concentrate“ vs. „Waschmittelkonzentrat“
    
    \item durch den Einfluss der Werbesprache auf Produkten und Schildern, bei der – unter Missachtung der deutschen Rechtschreibung – Wörter getrennt geschrieben werden, vermutlich um die schnelle Lesbarkeit in Sekundenbruchteilen durch potenzielle Käufer zu erleichtern. Falsch ist es trotzdem, auch wenn man Dinge wie „Waschmittel Konzentrat“, „Herren Schneiderei“ etc. allzu oft sieht.
\end{itemize}

Eine besondere Problematik speziell in der Wirtschaftsinformatik ist, dass es häufig Komposita aus englischen und deutschen Wörtern gibt. Hier kann man mit Bindestrichen arbeiten, wie z. B. „das Windows-Betriebssystem“. Zu beachten ist, dass immer alle Wortbestandteile des Kompositums mit Bindestrichen versehen werden müssen, auch wenn es einzeln nicht der Fall ist: „in Microsoft Windows“, aber „das Microsoft-Windows-Betriebssystem“. Rein englische Komposita können in der englischen Originalschreibweise verwendet werden, wie z. B. „program execution platform“, Großschreibung ist möglich: „Program Execution Platform“, als Kompositum mit dem deutschen Wort „Plattform“ wäre es aber wieder „Program-Execution-Plattform“.

\subsection{Kommasetzung}

Die Kommasetzung ist häufig fehlerhaft oder es werden Kommata „nach Gefühl“ gesetzt. Zugegebenermaßen sind die Kommaregeln im Deutschen komplex, für Fachtexte kommt man aber meistens mit einigen wenigen Regeln aus. Hier ein Auszug, wann Kommata zu setzen sind:

\begin{itemize}
    \item Relativsatz: \\
    „Wir verwenden eine Programmiersprache, die sehr leicht zu erlernen ist.“ Das Wort „die“ bezieht sich auf die Programmiersprache, der nachfolgende Nebensatz muss daher mit einem Komma abgetrennt werden.

    \item Konditionalsatz: \\
    „Wir verwenden diese Programmiersprache, weil sie sehr leicht zu erlernen ist.“ Das Wort „weil“ leitet einen Nebensatz ein, der der Begründung dient, und muss daher mit einem Komma abgetrennt werden.

    \item Einschub: \\
    „Wir verwenden diese Programmiersprache, übrigens sehr erfolgreich, in mehreren Projekten.“ Hier kann „übrigens sehr erfolgreich“ auch weggelassen werden und ist daher ein Einschub, der vorne und hinten durch ein Komma vom eigentlichen Satz abgetrennt wird. Ohne den Einschub hätte der Satz kein Komma.
\end{itemize}

Umgekehrt werden Kommata häufig fälschlich gesetzt, eine weit verbreitete Form ist das sog. „Vorfeldkomma“, mit dem eine kurze oder längere Phrase am Satzanfang abgetrennt wird, obwohl dies nicht korrekt ist. Beispiele für solche falschen Kommata sind:

\begin{itemize}
    \item „Das von Herrn Meier in der Schulung gezeigte Anwendungssystem, (falsches Komma!) wird auch in der Praxis häufig verwendet.“
    
    \item „Aufgrund des Sprachgefühls, (falsches Komma!) kommt hier ein Komma hin.“\footnote{Berg, Kristian, 2022, Aufgrund des Sprachgefühls kommt hier ein Komma hin \urlstyle{same}\mbox{\url{https://derzwiebel.wordpress.com/2022/02/22/vorfeldkomma/}}}
    
    \item „Neben dem vielfältigen Sportangebot, bietet die Isla Cozumel auch zahlreiche kulturelle Attraktionen“\footnote{Berg, Kristian, 2022, Aufgrund des Sprachgefühls kommt hier ein Komma hin \urlstyle{same}\mbox{\url{https://derzwiebel.wordpress.com/2022/02/22/vorfeldkomma/}}}
    
    \item „Studien von Pienemann (1998) zufolge, schwäche der Unterricht in einer zu hohen Profilstufe […] die LernerInnen rückwirkend.“\footnote{Berg, Kristian, 2022, Aufgrund des Sprachgefühls kommt hier ein Komma hin \urlstyle{same}\mbox{\url{https://derzwiebel.wordpress.com/2022/02/22/vorfeldkomma/}}}
\end{itemize}

Schauen Sie gerne in dem in der Fußnote verlinkten Dokument die Details nach, wann und warum Vorfeldkommata falsch sind und warum sie trotzdem häufig zu sehen sind (leider mittlerweile auch in redaktionellen Texten).

\subsection{Verwendung von Leerzeichen}

Dies ist ein Aspekt der Typografie, der sich insbesondere beim „Maschinenschreiben“ ausprägt und daher bei der überwiegend handschriftlichen Textproduktion in der Schule kaum berücksichtigt wird. Einige häufig auftretende Aspekte sollen beleuchtet werden:

\begin{itemize}
    \item Vor- und nach Schrägstrichen werden keine Leerzeichen gesetzt. Richtig ist also „schwarz/weiß“ und falsch wären „schwarz / weiß“ oder „schwarz /weiß“.

    \item Im Deutschen werden Leerzeichen vor und hinter Gedankenstriche gesetzt. Richtig ist also „Der Vogel – er war besonders schön – zwitscherte den ganzen Morgen“. Achtung, im Englischen verwendet man an dieser Stelle keine Leerzeichen, also nicht verwirren lassen. Generell sollte man sich überlegen, inwieweit Gedankenstriche nicht durch Kommata ersetzt werden sollten.

    \item Abkürzungen, die einzelne Wörter abkürzen, werden korrekt mit Leerzeichen gesetzt. Das häufigste Beispiel, das man auch am häufigsten falsch sieht, ist die Abkürzung für „zum Beispiel“. Hier muss korrekt „z.\,B.“ gesetzt werden, „z.B.“ ist falsch. Achten Sie einmal darauf, dass es in lektorierten Texten meist korrekt geschrieben wird, ansonsten überwiegend falsch. Da es manche kaum glauben wollen, dass es so ist, sei hier auf einige Quellen verwiesen:
    
    \begin{itemize}
        \item \verb|https://www.duden.de/rechtschreibung/z__B_ |
        \item \verb|https://www.din-5008-richtlinien.de/startseite/abkuerzungen/ |
        \item \verb|https://de.wikipedia.org/wiki/Geschütztes_Leerzeichen |
    \end{itemize}

    \item Weitere Abkürzungen, für die das gilt, sind „u.\,a.“, „d.\,h.“ etc. Bei „usw.“ sowie „etc.“ gilt eine Ausnahme. Damit Word keinen sehr unelegant aussehenden Zeilenumbruch an den Leerzeichen durchführt, sollte man statt eines normalen Leerzeichens bei solchen Abkürzungen immer ein „geschütztes Leerzeichen“ (\verb|\,|) eingeben, was nicht umgebrochen wird. Auch wird der Abstand zwischen den einzelnen Teilen der Abkürzung von Word dann etwas verringert.

    \item Das geschützte Leerzeichen sollte man an allen Stellen verwenden, wo man keinen Zeilenumbruch möchte, also auch z.\,B. bei „30.\,Juni“, „Seite\,15“, „Dr.\,Sommer“.

    \item Geschützte Leerzeichen werden auch zwischen Zahlenangabe und Maßeinheit gesetzt, also z.\,B. „15\,°C“, „100\,mm“, „180\,km/h“ etc.
\end{itemize}
